\section{完整约束}\label{sec:holonomic-constraints}

许多力学系统涉及位置受约束的粒子或粒子系统 (例如单摆, 在给定曲面上运动的粒子, 或用无质量杆连接起来的刚性粒子系统). 为处理这些情形, 引入以下定义.

\begin{RiemannDefinition}\label{def:holonomic-constraint}
力学系统 $(M,\langle\cdot,\cdot\rangle,\mathcal{F})$ 上的一个\textbf{完整约束}是满足 $\dim N<\dim M$ 的子流形 $N\subset M$. 若对所有 $t\in I$ 都有 $c(t)\in N$, 则称曲线 $c:I\subset\mathbb{R}\to M$ \textbf{与 $N$ 相容}.
\end{RiemannDefinition}

\begin{RiemannExample}\label{ex:holonomic-constraints-of-pendulum-surface-rigid-body}
\begin{enumerate}[label=(\arabic*),ref=\arabic*,leftmargin=*]
\item\label{ex:simple-pendulum-constraint} 一个质量为 $m>0$ 的粒子在 $\mathbb{R}^2$ 中运动, 受恒定重力加速度 $g$ 作用. 它由力学系统 $(\mathbb{R}^2,\langle\!\langle\cdot,\cdot\rangle\!\rangle,-mg\,\mathrm{d}y)$ 模拟, 其中
\[
  \langle\!\langle v,w\rangle\!\rangle:=m\langle v,w\rangle
\]
($\langle\cdot,\cdot\rangle$ 是 $\mathbb{R}^2$ 上的 Euclidean 内积). 用一根长度为 $l>0$ 的理想无质量杆将粒子与一个固定支点相连, 得到一个\textbf{单摆} (参见 Fig.~\ref{fig:simple-pendulum}). 假设支点是原点, 则这对应于完整约束
\[
  N=\{(x,y)\in\mathbb{R}^2\mid x^2+y^2=l^2\}
\]
($N$ 与 $S^1$ 微分同胚).
\item\label{ex:particle-on-surface-constraint} 类似地, 一个质量为 $m>0$ 的粒子在 $\mathbb{R}^3$ 中运动, 受恒定重力加速度 $g$ 作用. 它由力学系统 $(\mathbb{R}^3,\langle\!\langle\cdot,\cdot\rangle\!\rangle,-mg\,\mathrm{d}z)$ 模拟, 其中
\[
  \langle\!\langle v,w\rangle\!\rangle:=m\langle v,w\rangle
\]
($\langle\cdot,\cdot\rangle$ 是 $\mathbb{R}^3$ 上的 Euclidean 内积). 要求粒子在方程为 $z=f(x,y)$ 的曲面上运动, 得到完整约束
\[
  N=\{(x,y,z)\in\mathbb{R}^3\mid z=f(x,y)\}.
\]
\item\label{ex:rigid-body-distance-constraints} 质量分别为 $m_1,\ldots,m_k$ 的 $k$ 个粒子在 $\mathbb{R}^3$ 中自由运动. 这一系统由力学系统 $(\mathbb{R}^{3k},\langle\!\langle\cdot,\cdot\rangle\!\rangle,0)$ 模拟, 其中
\[
  \langle\!\langle(v_1,\ldots,v_k),(w_1,\ldots,w_k)\rangle\!\rangle
  :=\sum_{i=1}^k m_i\langle v_i,w_i\rangle
\]
($\langle\cdot,\cdot\rangle$ 是 $\mathbb{R}^3$ 上的 Euclidean 内积). 用理想无质量杆连接所有粒子得到\textbf{刚体}, 对应于完整约束
\[
  N=\bigl\{(x_1,\ldots,x_k)\in\mathbb{R}^{3k}\mid
    \lVert x_i-x_j\rVert=d_{ij}\ \text{对 }1\leq i<j\leq k\bigr\}.
\]
若至少有三个粒子不共线, 容易看出 $N$ 与 $\mathbb{R}^3\times O(3)$ 微分同胚.
\end{enumerate}
\end{RiemannExample}
\RGfigure[.46\textwidth]{simple-pendulum}{单摆}

要使粒子始终满足完整约束, 需要一个额外的外力 (在上述例子中由杆或曲面提供).

\begin{RiemannDefinition}\label{def:reaction-force-for-holonomic-constraint}
带完整约束的力学系统 $(M,\langle\cdot,\cdot\rangle,\mathcal{F},N)$ 上的\textbf{约束反力}是映射 $\mathcal{R}:TN\to T^*M$, 它对所有 $p\in N$ 满足 $\mathcal{R}(T_pN)\subset T_p^*M$, 且对每个 $v\in TN$, \textbf{广义 Newton 方程}
\[
  \mu\left(\frac{D\dot c}{\mathrm{d}t}\right)
  =(\mathcal{F}+\mathcal{R})(\dot c)
\]
有一个满足初始条件 $\dot c(0)=v$ 的解 $c:I\subset\mathbb{R}\to N$.
\end{RiemannDefinition}

对任意完整约束, 一般有无穷多种约束反力可供选择. 以下定义给出选择约束反力的一条特别有用的判据.

\begin{RiemannDefinition}\label{def:perfect-reaction-force}
若带完整约束的力学系统 $(M,\langle\cdot,\cdot\rangle,\mathcal{F},N)$ 的约束反力对所有 $v\in T_pN$ 和 $p\in N$ 满足
\[
  \mu^{-1}(\mathcal{R}(v))\in(T_pN)^\perp,
\]
则称其为\textbf{理想约束反力}, 或称它满足 \textbf{d’Alembert 原理}.
\end{RiemannDefinition}

\begin{RiemannRemark}\label{rem:perfect-reaction-does-no-work}
广义 Newton 方程的解的动能变化为
\begin{align*}
  \frac{\mathrm{d}K}{\mathrm{d}t}
  &=\left\langle\frac{D\dot c}{\mathrm{d}t},\dot c\right\rangle
  =\mathcal{F}(\dot c)(\dot c)+\mathcal{R}(\dot c)(\dot c)\\
  &=\mathcal{F}(\dot c)(\dot c)+\bigl\langle\mu^{-1}(\mathcal{R}(\dot c)),\dot c\bigr\rangle.
\end{align*}
因此, 约束反力是理想的, 当且仅当沿任何与约束相容的运动, 它既不产生能量也不耗散能量.
\end{RiemannRemark}

\begin{RiemannExample}\label{ex:perfect-reactions-in-three-holonomic-systems}
在上述各例中, 要求约束反力为理想约束反力分别等价于下列假设.
\begin{enumerate}[label=(\arabic*),leftmargin=*]
\item \textbf{单摆:} 杆传递的力完全沿径向 (即没有阻尼);
\item \textbf{曲面上的粒子:} 曲面施加的力与曲面正交 (即曲面无摩擦);
\item \textbf{刚体:} 内聚力不耗散能量.
\end{enumerate}
\end{RiemannExample}

下述结果建立理想约束反力的存在性和唯一性.

\begin{RiemannTheorem}\label{thm:existence-and-uniqueness-of-perfect-reaction}
给定任意带完整约束的力学系统 $(M,\langle\cdot,\cdot\rangle,\mathcal{F},N)$, 存在唯一满足 d’Alembert 原理的约束反力 $\mathcal{R}:TN\to T^*M$. 广义 Newton 方程
\[
  \mu\left(\frac{D\dot c}{\mathrm{d}t}\right)
  =(\mathcal{F}+\mathcal{R})(\dot c)
\]
的解恰是力学系统 $(N,\langle\!\langle\cdot,\cdot\rangle\!\rangle,\mathcal{F}_N)$ 的运动, 其中 $\langle\!\langle\cdot,\cdot\rangle\!\rangle$ 是由 $\langle\cdot,\cdot\rangle$ 在 $N$ 上诱导的度量, $\mathcal{F}_N$ 是 $\mathcal{F}$ 到 $N$ 的限制. 特别地, 若 $\mathcal{F}=-\mathrm{d}U$ 是保守力, 则 $\mathcal{F}_N=-\mathrm{d}(U|_N)$.
\end{RiemannTheorem}

\begin{RiemannProof}
回顾第~\ref{chap:curvature}~章第~\ref{sec:isometric-immersions}~节, 若 $\widetilde\nabla$ 是 $(M,\langle\cdot,\cdot\rangle)$ 的 Levi–Civita 联络, $\nabla$ 是 $(N,\langle\!\langle\cdot,\cdot\rangle\!\rangle)$ 的 Levi–Civita 联络, 则
\[
  \nabla_XY=(\widetilde\nabla_{\widetilde X}\widetilde Y)^\top
\]
对所有 $X,Y\in\mathfrak{X}(N)$ 成立, 其中 $\widetilde X,\widetilde Y$ 是 $X,Y$ 到 $\mathfrak{X}(M)$ 的任意延拓 (照常, $v=v^\top+v^\perp$ 表示对每个 $p\in N$ 由分解 $T_pM=T_pN\oplus(T_pN)^\perp$ 得到的唯一分解). 此外, $N$ 的第二基本形式
\[
  B(X,Y)=\widetilde\nabla_{\widetilde X}\widetilde Y-\nabla_XY
        =(\widetilde\nabla_{\widetilde X}\widetilde Y)^\perp
\]
是良定义的, 且对所有 $p\in N$, $B(X,Y)_p\in(T_pN)^\perp$ 是 $X_p,Y_p$ 的对称双线性函数.

假设理想约束反力 $\mathcal{R}$ 存在, 则广义 Newton 方程的解满足
\[
  \widetilde\nabla_{\dot c}\dot c=\mu^{-1}(\mathcal{F}(\dot c))+\mu^{-1}(\mathcal{R}(\dot c)).
\]
由于按假设 $\mu^{-1}\mathcal{R}$ 垂直于 $N$, 该方程沿 $N$ 的切向分量给出
\[
  \nabla_{\dot c}\dot c=\mu_N^{-1}(\mathcal{F}_N(\dot c))
\]
(其中 $\mu_N:TN\to T^*N$ 是 $N$ 上的质量算子), 因为对任意 $v\in TN$ 有
\begin{align*}
  \left\langle\!\left\langle(\mu^{-1}(\mathcal{F}(\dot c)))^\top,v\right\rangle\!\right\rangle
  &=\left\langle\mu^{-1}(\mathcal{F}(\dot c)),v\right\rangle=\mathcal{F}(\dot c)(v)\\
  &=\mathcal{F}_N(\dot c)(v)=\left\langle\!\left\langle\mu_N^{-1}(\mathcal{F}_N(\dot c)),v\right\rangle\!\right\rangle.
\end{align*}
因此 $c$ 是 $(N,\langle\!\langle\cdot,\cdot\rangle\!\rangle,\mathcal{F}_N)$ 的一个运动.

另一方面, 广义 Newton 方程垂直于 $N$ 的分量给出
\[
  B(\dot c,\dot c)=(\mu^{-1}(\mathcal{F}(\dot c)))^\perp+\mu^{-1}(\mathcal{R}(\dot c)).
\]
所以, 若 $\mathcal{R}$ 存在, 则对所有 $v\in TN$ 必须满足
\begin{equation}\label{eq:perfect-holonomic-reaction-force}
  \mathcal{R}(v)=\mu(B(v,v))-\mu\left[(\mu^{-1}(\mathcal{F}(v)))^\perp\right].
\end{equation}
这就证明了唯一性.

为证明存在性, 由 \eqref{eq:perfect-holonomic-reaction-force} \textbf{定义} $\mathcal{R}$. 这当然保证对所有 $v\in T_pN$ 和 $p\in N$, 有 $\mu^{-1}(\mathcal{R}(v))\in(T_pN)^\perp$. 给定 $v\in TN$, 设 $c:I\subset\mathbb{R}\to N$ 是力学系统 $(N,\langle\!\langle\cdot,\cdot\rangle\!\rangle,\mathcal{F}_N)$ 中以 $v$ 为初始条件的运动. 则
\begin{align*}
  \widetilde\nabla_{\dot c}\dot c
  &=\nabla_{\dot c}\dot c+B(\dot c,\dot c)\\
  &=\mu_N^{-1}(\mathcal{F}_N(\dot c))
     +(\mu^{-1}(\mathcal{F}(\dot c)))^\perp+\mu^{-1}(\mathcal{R}(\dot c))\\
  &=(\mu^{-1}(\mathcal{F}(\dot c)))^\top
     +(\mu^{-1}(\mathcal{F}(\dot c)))^\perp+\mu^{-1}(\mathcal{R}(\dot c))\\
  &=\mu^{-1}(\mathcal{F}(\dot c))+\mu^{-1}(\mathcal{R}(\dot c)).
\end{align*}
\end{RiemannProof}

\begin{RiemannExample}\label{ex:simple-pendulum-equation-of-motion}
为写出带理想约束反力的单摆的运动方程, 用映射 $\varphi:(-\pi,\pi)\to\mathbb{R}^2$ 参数化完整约束 $N$, 该映射定义为
\[
  \varphi(\theta)=(l\sin\theta,-l\cos\theta)
\]
(所以 $\theta=0$ 标记稳定平衡位置, 参见 Fig.~\ref{fig:simple-pendulum}). 有
\[
  \frac{\mathrm{d}}{\mathrm{d}\theta}
  =\frac{\mathrm{d}x}{\mathrm{d}\theta}\frac{\partial}{\partial x}
   +\frac{\mathrm{d}y}{\mathrm{d}\theta}\frac{\partial}{\partial y}
  =l\cos\theta\,\frac{\partial}{\partial x}+l\sin\theta\,\frac{\partial}{\partial y},
\]
所以单摆的动能为
\begin{align*}
  K\left(v\frac{\mathrm{d}}{\mathrm{d}\theta}\right)
  &=\frac12m\left\langle
      vl\cos\theta\,\frac{\partial}{\partial x}+vl\sin\theta\,\frac{\partial}{\partial y},\
      vl\cos\theta\,\frac{\partial}{\partial x}+vl\sin\theta\,\frac{\partial}{\partial y}
      \right\rangle\\
  &=\frac12ml^2v^2.
\end{align*}
另一方面, 势能由下式给出:
\[
  U(x,y)=mgy,
\]
所以它到 $N$ 的限制的局部表达式为
\[
  U(\theta)=-mgl\cos\theta.
\]
因而运动方程为
\begin{align*}
  &\frac{\mathrm{d}}{\mathrm{d}t}\left(\frac{\partial K}{\partial v}(\theta,\dot\theta)\right)
   -\frac{\partial K}{\partial\theta}(\theta,\dot\theta)=-\frac{\partial U}{\partial\theta}(\theta)\\
  &\quad\Leftrightarrow\frac{\mathrm{d}}{\mathrm{d}t}(ml^2\dot\theta)=-mgl\sin\theta\\
  &\quad\Leftrightarrow\ddot\theta=-\frac gl\sin\theta.
\end{align*}
注意, 无需求出约束反力.
\end{RiemannExample}

\begin{RiemannExercise}\label{ex:holonomic-constraints}
\begin{enumerate}[label=(\arabic*),ref=\arabic*,leftmargin=*]
\item\label{ex:spherical-pendulum-motion} 利用球坐标写出长度为 $l$ 的\textbf{球面摆}的运动方程. 它是质量为 $m>0$ 的粒子在 $\mathbb{R}^3$ 中的运动, 受恒定重力加速度 $g$ 和完整约束
\[
  N=\{(x,y,z)\in\mathbb{R}^3\mid x^2+y^2+z^2=l^2\}.
\]
$N$ 上哪些纬线是粒子的可能轨迹?
\item\label{ex:particle-on-frictionless-surface-of-revolution} 写出粒子在恒定重力加速度 $g$ 作用下沿方程为 $z=f(r)$ (其中 $r=\sqrt{x^2+y^2}$) 的无摩擦旋转曲面运动的方程.
\item\label{ex:free-dumbbell-motion} 写出并求解自由哑铃 (即由长度为 $l$ 的无质量杆连接的质量分别为 $m_1,m_2$ 的两个粒子组成的系统) 在下列空间内运动的方程:
\begin{enumerate}[label=(\alph*),leftmargin=*]
\item $\mathbb{R}^2$;
\item $\mathbb{R}^3$.
\end{enumerate}
{\small (\textbf{提示:} 使用\textbf{质心}的坐标, 即沿杆与 $m_1$ 相距 $\frac{m_2}{m_1+m_2}l$ 的点的坐标).}
\item\label{ex:double-pendulum-motion} 长度为 $l_1,l_2$ 的\textbf{双摆}是由质量分别为 $m_1,m_2$ 的两个粒子确定的力学系统. 粒子在 $\mathbb{R}^2$ 中运动, 受恒定重力加速度 $g$ 和完整约束
\[
  N=\{(x_1,x_2)\in\mathbb{R}^4\mid\lVert x_1\rVert=l_1\ \text{且}\ \lVert x_1-x_2\rVert=l_2\}.
\]
($N$ 与 2 维环面 $T^2$ 微分同胚.)
\begin{enumerate}[label=(\alph*),leftmargin=*]
\item 用以下参数化 $\phi:(-\pi,\pi)\times(-\pi,\pi)\to N$ 写出双摆的运动方程:
\[
  \phi(\theta,\varphi)=
  (l_1\sin\theta,-l_1\cos\theta,l_1\sin\theta+l_2\sin\varphi,-l_1\cos\theta-l_2\cos\varphi)
\]
(参见 Fig.~\ref{fig:double-pendulum}).
\item 在 $\theta=\varphi=0$ 附近将运动方程线性化. 寻找线性化方程满足 $\varphi=k\theta$ 的解, 其中 $k\in\mathbb{R}$ 是常数 (\textbf{简正模}). 由此产生的振荡的周期分别是多少?
\end{enumerate}
\end{enumerate}
\end{RiemannExercise}
\RGfigure[.43\textwidth]{double-pendulum}{双摆}
